%% expect: native native fallback fallback native native fallback native native fallback native fallback native native
%% reason 3: a tip redefined
%% reason 4: a tip redefined
%% reason 7: an arrow tip set by >
%% reason 10: an arrow tip set by >
%% reason 12: every loop
\documentclass{article}
\usepackage{tikz-cd}
\usetikzlibrary{arrows.meta}
\newbox\unitbox
\makeatletter
\newcommand\dims[1]{\typeout{[[#1 \the\wd\unitbox\space\the\ht\unitbox\space\the\dp\unitbox]]}}
\makeatother
\begin{document}

% A tip of tikz-cd's redefined reaches the arrows that draw it, and those alone.
{\tikzcdset{tikzcd left hook/.tip={Bar[width=8pt]}}
\setbox\unitbox=\hbox{\begin{tikzcd} A \arrow[r] & B \end{tikzcd}}\dims{1}\box\unitbox
\setbox\unitbox=\hbox{\begin{tikzcd} A \arrow[r,two heads,dashed] & B \end{tikzcd}}\dims{2}\box\unitbox
\setbox\unitbox=\hbox{\begin{tikzcd} A \arrow[r,hook'] & B \end{tikzcd}}\dims{3}\box\unitbox
% The right hook is drawn with the left one.
\setbox\unitbox=\hbox{\begin{tikzcd} A \arrow[r,hook] & B \end{tikzcd}}\dims{4}\box\unitbox
}
\setbox\unitbox=\hbox{\begin{tikzcd} A \arrow[r,hook] & B \end{tikzcd}}\dims{5}\box\unitbox
% > reaches the arrows of /tikz, and what the library knows of it it takes.
{\tikzset{>=stealth}
\setbox\unitbox=\hbox{\begin{tikzcd} A \arrow[r,->] \arrow[d] & B \\ C \end{tikzcd}}\dims{6}\box\unitbox
}
{\tikzset{>={Latex[length=5pt]}}
\setbox\unitbox=\hbox{\begin{tikzcd} A \arrow[r,->] & B \end{tikzcd}}\dims{7}\box\unitbox
\setbox\unitbox=\hbox{\begin{tikzcd} A \arrow[r] & B \end{tikzcd}}\dims{8}\box\unitbox
}
{\tikzset{>=latex}\tikzcdset{arrow style=tikz}
\setbox\unitbox=\hbox{\begin{tikzcd} A \arrow[r,two heads] & B \end{tikzcd}}\dims{9}\box\unitbox
\tikzset{>={Latex[length=5pt]}}
\setbox\unitbox=\hbox{\begin{tikzcd} A \arrow[r] & B \end{tikzcd}}\dims{10}\box\unitbox
}
% every loop is run by a loop's path.
{\tikzset{every loop/.style={red}}
\setbox\unitbox=\hbox{\begin{tikzcd} A \arrow[r] & B \end{tikzcd}}\dims{11}\box\unitbox
\setbox\unitbox=\hbox{\begin{tikzcd} A \arrow[loop right] & B \end{tikzcd}}\dims{12}\box\unitbox
}
\setbox\unitbox=\hbox{\begin{tikzcd} A \arrow[loop right] & B \end{tikzcd}}\dims{13}\box\unitbox
\setbox\unitbox=\hbox{\begin{tikzcd} A \arrow[r] & B \end{tikzcd}}\dims{14}\box\unitbox

\end{document}
